\documentclass[10pt,a4paper]{amsart}
\usepackage[margin=19mm]{geometry}
\usepackage{amsmath,amssymb,mathtools}
\usepackage[scr=rsfs,cal=euler]{mathalfa}
\usepackage{xfrac}
\usepackage[unicode,hidelinks,bookmarks=false]{hyperref}
\newcommand{\ie}{i.e.\ }
\newcommand{\cf}{cf.\ }
\newcommand{\resp}{respectively\ }
\newcommand{\Subsection}{\subsection}
\providecommand{\nobreakdash}{\nobreak\hbox{-}}
\setlength{\emergencystretch}{2em}
\multlinegap0pt
\usepackage{amssymb}
\usepackage{dsfont}
\usepackage{mathtools}
\usepackage{caption}
\usepackage[shortlabels]{enumitem}
\usepackage{tikz}
\usepackage[all]{xy}
\usepackage{xcolor}
\newtheorem{proposition}{Proposition}
\newcommand\bC{\mathbb{C}}
\newcommand{\ob}[1]{\mkern 1.5mu\overline{\mkern-1.5mu#1\mkern-1.5mu}\mkern 1.5mu}
\title{A canonical generator for congruence ideals of Hida families}
\author{Alexandre Maksoud}
\date{Source: arXiv:2312.16706v1 (CC BY 4.0)}
\begin{document}
\maketitle

\noindent\emph{Source and licence.} A canonical generator for congruence ideals of Hida families. Version arXiv:2312.16706v1 (CC BY 4.0), distributed by arXiv under the Creative Commons Attribution 4.0 International licence, http://creativecommons.org/licenses/by/4.0/, as recorded in the arXiv licence metadata for this exact version and shown on the abstract page https://arxiv.org/abs/2312.16706v1. Full licence notice: \texttt{licenses/arxiv-2312.16706-CC-BY-4.0.txt}.\par\smallskip

\noindent\textit{Editorial scope.} Counted unit: the article's proposition prop:calcul\_poincare\_pairing (source0.tex lines 551--557). The article's complete original proof of it is delivered with it (source0.tex lines 558--577, one \texttt{proof} environment of 20 lines). No mathematics is written, repaired or re-derived: the statement and the proof are the article's own lines, in the article's own notation, and no retained line is substituted or re-flowed.  Retained uncounted as the source-local context the proof needs: lines 500--502.  Nothing was omitted from any retained span, so the package carries no omission record.  No retained line contains a citation, so the wrapper carries no bibliography and no bibliography entry.  No retained line contains a figure environment, an image-inclusion command or drawing code, so no image is delivered and the package declares no asset dependency.  Editorial formatting only, in addition: an AMS \texttt{amsart} preamble that restates the article's own declarations and macros used by the retained lines, namely the environments \texttt{proposition} on separate counters, together with the standard packages those lines need.  The retained environments print in this document as Proposition 1 in order of appearance, whereas the article prints the same items with the numbering of its own full document.  The complete native source of the article as distributed in the arXiv e-print for this exact version is bundled unchanged as \texttt{sources/152142/source0.tex}.
\begin{equation}\label{eq:action_slash_cohomology}
	\delta(f)|[\Gamma s \Gamma]=\delta(f|[\Gamma s \Gamma]) \qquad \mbox{and}\qquad \delta^c(f)|[\Gamma s \Gamma]=\delta^c(f|[\Gamma \iota s\iota \Gamma]).
\end{equation}

\begin{proposition}\label{prop:calcul_poincare_pairing}
	Let $f,g\in S_k(M)$. We have 
	\begin{align*}
		(\delta(f),\delta^c(g))_\pi&= (2i)^{k-1}\langle f,g^c\rangle_\Gamma \\
		(\delta^+(f),\delta^-(g)|W_M)_\pi&=-2(2i)^{k-1}\langle f,g^c|W_M\rangle_\Gamma.
	\end{align*}
\end{proposition}

\begin{proof}
	We start with the first equality. Over $\bC$, the cap-product with the fundamental class of $X(\bC)$ is given by the integration over $X(\bC)$ of cohomology classes seen as $2$-differential forms. Therefore,
	\begin{align*}
		(\delta(f),\delta^c(g))_\pi &= - \int_{X(\bC)}   f(z)\ob{g^c(z)} [(X-zY)^{k-2},(X-\ob{z}Y)^{k-2}] dz\wedge d\ob{z} \\
		&= 2i \int_{X(\bC)} f(z)\ob{g^c(z)}(z-\ob{z})^{k-2} dxdy \\
		&=(2i)^{k-1}\int_{X(\bC)} f(z)\ob{g^c(z)}y^{k-2} dxdy.
	\end{align*}
This proves the first formula. A similar computation as above yields \small
\begin{align*}
	(\delta(h),\delta(h'))_\pi&=0=(\delta^c(h),\delta^c(h'))_\pi, \\  (\delta(h),\delta^c(h'))_\pi&=(-1)^{k-1}(\delta^c(h'),\delta(h))_\pi,
\end{align*}
\normalsize for any $h,h'\in S_k(M)$, and from \eqref{eq:action_slash_cohomology} we find $\delta^-(g)|W_M=\delta(g|W_M)+(-1)^{k-1}\delta^c(g|W_M)$. Hence, by linearity,
\begin{align*}
	(\delta^+(f),\delta^-(g)|W_M)_\pi&= (-1)^{k-1} \left( (\delta(f),\delta^c(g|W_M))_\pi + (\delta(g|W_M),\delta^c(f))_\pi \right) \\
									&=(-2i)^{k-1}\left(\langle f,(g|W_M)^c\rangle_\Gamma+\langle g|W_M,f^c\rangle_\Gamma\right) \\
									&= 2(-2i)^{k-1}\langle f,(g|W_M)^c\rangle_\Gamma \\
									&=-2(2i)^{k-1}\langle f,g^c|W_M\rangle_\Gamma,
\end{align*}
since $(g|W_M)^c=(-1)^kg^c|W_M$. This proves the second formula.
\end{proof}

\end{document}
